% Version 0.27.  Actual arithmetic data and the complete level comparison.
% Pinned Lean source: 7cbf7a0c32fca1badf05931879a37c3f8a2b03b4.
\section{The actual symplectic arithmetic data}
\label{sec:symplectic-arithmetic-data}

\status{AX1--AX10 close the arithmetic construction behind G1.
They prove the literal \texttt{dataEx}, all of \texttt{facts},
\texttt{levOf}, the ordinary specialization theorem \texttt{spPtL},
and the level, duality and safety assertions on the actual points.
The final admissible class uses the explicitly retained
\texttt{SymplecticFourImageDense} hypothesis to supply its common
automorphic embedding, \texttt{Fin0} and \texttt{NoEnd}.  None of those
three assertions is inferred from the character construction.}

\subsection{The original objects and their conditions}

Fix the source algebraic closure $\Omega/\mathbb Q$ and put
$G=\operatorname{Gal}(\Omega/\mathbb Q)$.  Let $\ell$ be an odd prime,
let $O\subset\overline{\mathbb Q}_\ell$ be the original coefficient
DVR, finite over $\mathbb Z_\ell$, and let $k=O/\mathfrak m_O$.
Write $c:G\to\mathbb Z_\ell^\times\subset O^\times$ for the
cyclotomic character, $\omega$ for its Teichm\"uller part, and a bar
for reduction to $k$.  A Dirichlet character is evaluated on $G$
through the cyclotomic action at its specified level.  Characters of
different levels are multiplied after inflation to a common multiple;
all displayed identities involving them mean identities on units.

Keep the host of HV1: a primitive character $\chi$ of conductor $q$
and order $o$, an integer $s\ge2$, and the original prime pair $q_1,d$.
Put
\[
 N=\operatorname{lcm}(2^{s+2},q)q_1d,
 \qquad H\le(\mathbb Z/N)^\times
 \tag{AX.1}
\]
with exactly the four defining conditions in HV1.  A host character
$\psi$ is trivial on $H$.  Membership in $h.X(O,e)$ also requires
$\psi(-1)=(-1)^e$.  Throughout, $\operatorname{clog}_2(x)$ is the
least $n\ge0$ with $x\le2^n$ for the positive integers used below.

For a finite-order character, let $\operatorname{comp}_p$ denote its
CRT component supported at $p$.  Split
\[
 \psi=\psi_2\theta\kappa,
 \quad \psi_2=\operatorname{comp}_2\psi,
 \quad \kappa=\psi_{\ell\text{-primary}},
 \quad \theta=\psi_{\ell'}(\operatorname{comp}_2\psi_{\ell'})^{-1}.
 \tag{AX.2}
\]
Also put $\psi_P=\psi_{\ell'}
(\operatorname{comp}_\ell\psi_{\ell'})^{-1}$,
$c_2=v_2(\operatorname{cond}\psi_2)$,
$f_0=\operatorname{ord}_{\operatorname{lcm}(4,\operatorname{ord}\theta)}(\ell)$,
and $n_0=v_2(\ell^{f_0}-1)$.  These are the literal source
\texttt{psi2}, \texttt{th1}, \texttt{kap}, \texttt{chiP},
\texttt{c2}, \texttt{f0} and \texttt{n0}.
The two prescribed dyadic levels are
\[
\begin{split}
 n_{\rm II}&=\operatorname{clog}_2\max
  \{\ell N4^{n_0},\ 2^{c_2+6}o\ell^3q^2\},\\
 n_{\rm III}&=\operatorname{clog}_2\max
  \{\ell q p_{\!1}4^{n_0},\ 8\ell^2\}.
\end{split}\tag{AX.3}
\]
Here $p_{\!1}=\mathtt{pOne}(\chi)$ is the fixed prime congruent to
$1$ modulo the odd part of $o$ and coprime to $2qp_K$.  It does not
vary with the host.  The exceptional set is the same
$\mathcal B_\chi$ as in HV1.  In the two cases,
\[
\begin{array}{ll}
\mathrm{II}:&e\ge2,\quad \operatorname{ord}(\psi_P)\text{ even},
 \quad[\ell\in\mathcal B_\chi\ \text{or}\ \ell-1\mid2e],\\
\mathrm{III}:&\ell\notin\mathcal B_\chi,\quad\ell-1\mid2e,
 \quad\operatorname{ord}(\psi)\text{ odd},\quad\psi\ne1.
\end{array}\tag{AX.4}
\]

A source \texttt{Raw} datum $D$ consists of $S,e,\psi,\xi,\omega$,
a complex conjugation $\mathrm{cc}$, a place $A_p$ and an element
$\sigma_p$ for every rational prime $p$, and $\phi,g_0,s_0$.
The same symbol $\psi$ will denote its associated Galois character.
The derived characters and four weights are
\[
 \begin{gathered}
 \Psi=c^e\psi,\quad \alpha=c\omega^{-1}\xi,
 \quad\beta=\Psi\alpha^{-1},\quad\vartheta=\alpha\Psi,\\
 (p_0,p_1,p_2,p_3)=(\vartheta,\Psi,\alpha,1),
 \quad(w_0,w_1,w_2,w_3)=(e+1,e,1,0).
 \end{gathered}
 \tag{AX.5}
\]
Here the subscripts on $p_i$ denote characters, not rational primes.
Let $r_0=\operatorname{diag}(p_i)$ and $\bar r_0=\overline{r_0}$.
The source local sets are
\[
 S_b=\{p\in S:p\ne\ell,
       \ \ell\mid|\psi(I_{A_p})|\},\qquad T=\{\ell\}\cup S_b.
 \tag{AX.6}
\]
The condition $\operatorname{UnrT}(z,0)$ says that $z(g)=0$ for
$g\in I_A\cap\ker\bar r_0$ whenever $A$ is away from $T$.
The condition $X\!\operatorname{Van}(\gamma)$ says that a locally
constant $k(\bar\gamma)$-cocycle with this property and locally a
coboundary on each $G_{A_v}$, $v\in T$, is a global coboundary.

For clarity, the full \texttt{Good} condition has the following
content.  The set $S$ contains $2,\ell$ and consists of primes;
$e\ge2$; the selected $A_p$ lie over $p$; and $\mathrm{cc}$ is a
complex conjugation.  The character $\psi$ has open kernel, is
unramified outside $S$, and $\Psi(\mathrm{cc})=1$.
One has $\omega^{\ell-1}=1$, $\bar\omega=\bar c$.
The character $\xi$ has open kernel and order prime to $\ell$,
is unramified outside $S$ and at all $v\in T$, and is odd.
The elements $\sigma_v$ are arithmetic Frobenius lifts for $v\in T$,
$c(\sigma_\ell)=1$, and
\[
 \overline{\xi\Psi^i}(\sigma_v)\ne1
       \quad(v\in T,\ i\in\mathbb Z).
 \tag{AX.7}
\]
Both $X\!\operatorname{Van}(\gamma)$ and
$X\!\operatorname{Van}(\gamma^{-1})$ hold for
$\gamma=\alpha,\beta,\vartheta$.  Finally,
\[
 \begin{gathered}
 \sigma_\ell^{\ell^{n!}}\longrightarrow\phi,
 \qquad\bar\alpha(g_0)\ne1,\quad\bar\Psi(g_0)\ne1,\\
 \qquad s_0\in I_{A_\ell}\cap\ker\psi\cap\ker\omega,
 \quad c(s_0)\text{ has infinite order}.
 \end{gathered}
 \tag{AX.8}
\]
This list includes every field of the source record; no stronger
goodness assumption will be added later.

\paragraph{Classical boundary.}
We use the elementary finite cyclotomic and local-field facts in
Milne, \href{https://www.jmilne.org/math/CourseNotes/ANT.pdf}
{\emph{Algebraic Number Theory}}, v3.08, Sections 6--8: the
cyclotomic Galois group, extension and conjugacy of places, the
decomposition-group comparison with the local completion, and the
wild/tame inertia description.  The exact consequences used below
are proved from these finite statements.  The norm on a finite
extension of $\mathbb Q_\ell$ is its unique extension, as in RV1.
The one commutative-algebra fact used in AX8 is that a finite algebra
over a Henselian local ring is a product of Henselian local rings
(\href{https://stacks.math.columbia.edu/tag/04GH}{Stacks, Tag 04GH});
a domain is therefore local.  Completeness implies Henselianity by
successive lifting of simple roots
(\href{https://stacks.math.columbia.edu/tag/04GM}{Tag 04GM}).
BR7 supplies the full
residual vanishing theorem, with its already recorded classical
inputs.  There is no automorphy, crystalline-period, image-density,
or rigidity input before AX10.

\subsection{Finite characters with the prescribed conductors}

\begin{proposition}[AX1: exact host arithmetic and the two cases]
\label{prop:ax-host}
For a host character the decomposition (AX.2) holds.  Its dyadic
component has conductor $2^{c_2}$ and order dividing $2^{c_2}$;
$\kappa$ has $\ell$-power order; and $\theta$ has odd conductor,
prime-to-$\ell$ order dividing $2o$.
In particular $f_0>0$, $4\mid\ell^{f_0}-1$,
$\theta^{\ell^{f_0}-1}=1$, and $\ell-1\mid\ell^{f_0}-1$.

Put $T_b=\{\ell\}\cup\{p\mid N:p\ne\ell,
\operatorname{comp}_p\kappa\ne1\}$.  Every $p\in T_b\setminus\{\ell\}$
is odd, divides $q$, and is congruent to $1$ modulo $\ell$.
Every $v\in T_b$ divides $\ell q$.
For any places chosen over the rational primes, the actual set
$T$ in (AX.6), with $D.\psi=\psi$ and prime-valued $S$, lies in $T_b$.

In Case II, $\psi_P\ne1$ and $n_{\rm II}\ge c_2+6$.
In Case III, under the original $h.\mathtt{Big}(t)$ and
$0<e<t$, one has
\[
 \begin{gathered}
 e\ge2,\quad\ell\nmid oN,\quad\ell\ne p_{\!1},
 \quad\psi_2=\kappa=1,\quad\theta=\psi_P=\psi,\\
 \quad\operatorname{cond}\psi\mid q,\quad T_b=\{\ell\},
 \quad\operatorname{ord}\psi\mid\operatorname{oddpart}(o),
 \quad n_{\rm III}\ge3.
 \end{gathered}
 \tag{AX.9}
\]
\end{proposition}

\begin{proof}
The finite unit group modulo $N$ is the product of its prime-power
unit groups.  Projection to a factor gives
$\operatorname{comp}_p$, so the components commute, their product
is the original character, and their orders divide the original
order.  Splitting a finite cyclic character image into its
$\ell$-primary and prime-to-$\ell$ factors commutes with every
component.  The group of units modulo a power of $2$ is a $2$-group;
thus $\operatorname{comp}_2\kappa=1$.  Removing the remaining dyadic
component gives (AX.2), odd conductor for $\theta$, and the asserted
dyadic order bound by the size of the unit group.

For each unit $u$ modulo $N$, the power $u^{2^{s+1}o}$ lies in $H$:
it is $1$ modulo $2^{s+2}$, is killed by $\chi$, and is a square
at the two auxiliary primes.  Hence
$\operatorname{ord}\psi\mid2^{s+1}o$.  More precisely, for an odd
prime $p\mid q$ the $p$-supported unit raised to $o$ lies in $H$,
so $\operatorname{ord}(\operatorname{comp}_p\psi)\mid o$.
For an odd prime $p\nmid q$ its square lies in $H$, so this
component has order dividing $2$.  Consequently every odd
component, and then their product after taking prime-to-$\ell$
parts, has order dividing $2o$.  The unit $\ell$ modulo
$\operatorname{lcm}(4,\operatorname{ord}\theta)$ has finite positive
order $f_0$, which gives its stated divisibilities.  The last one
also follows directly from $\ell^f-1=(\ell-1)(1+\cdots+\ell^{f-1})$.

A nontrivial $\ell$-primary component at an odd prime $p\ne\ell$
cannot occur at either auxiliary prime: there the character has
order dividing $2$.  Thus $p\mid q$.  On
$(\mathbb Z/p^a)^\times$ the kernel of reduction modulo $p$ is
a $p$-group, so an $\ell$-primary character factors through
$(\mathbb Z/p)^\times$ and its nontrivial order divides $p-1$.
This proves the assertions about $T_b$.  Inertia on roots of unity
of order prime to $p$ is trivial.  Therefore the restriction of
$\psi$ to $I_{A_p}$ equals the restriction of its $p$-component.
If the former image has order divisible by $\ell$, so does the
latter.  This is equivalent to
$\operatorname{comp}_p\kappa\ne1$, proving $T\subset T_b$ without
changing the selected places.

Even order makes $\psi_P$ nontrivial.  The second term in the
maximum defining $n_{\rm II}$ is at least $2^{c_2+6}$, so its
claimed lower bound follows.  For Case III, an odd-order character
has trivial dyadic component and trivial components at the two
auxiliary primes, hence conductor dividing $q$ and order dividing
the odd part of $o$.  Since $\ell\notin\mathcal B_\chi$, the
prime $\ell$ is at least $5$, does not divide $oq$, and differs
from $p_{\!1}$.  Also $\ell-1\mid2e$ gives
$\ell\le2e+1<2t+1<q_1,d$.  Thus $\ell\nmid N$ and $e\ge2$.
The $\ell$-primary and $\ell$-local components of $\psi$ vanish,
which proves all identifications in (AX.9).  Finally $8\ell^2$
in (AX.3) gives $n_{\rm III}\ge3$.

Only the fixed defining property of $p_{\!1}$ is used here.
Its existence in the stated class can even be seen without an
analytic progression theorem.  For $m>1$, choose an integer $a$
divisible by $m$ and any prescribed finite excluded set, with
$a\ge3$.  Any prime divisor $p$ of $\Phi_m(a)$ is outside that
set and does not divide $m$, since $\Phi_m(a)\equiv1\pmod a$.
The reduction of $X^m-1$ is separable at $p$, so a root of
$\Phi_m$ has order exactly $m$, and $p\equiv1\pmod m$.
The inequality $|\Phi_m(a)|\ge(a-1)^{\varphi(m)}>1$ supplies a
prime divisor.  For $m=1$ use the usual Euclidean construction.
Taking $m=\operatorname{oddpart}(o)$ and excluding primes dividing
$2qp_K$ proves that the least such prime has precisely the
properties used above.  The earlier fixed definitions of $p_K$
and $\mathcal B_\chi$ retain their meanings.
\end{proof}

\begin{proposition}[AX2: primitive dyadic characters and quantitative separation]
\label{prop:ax-dyadic}
If $n\ge3$ and $O$ contains all $2^n$-th roots of unity, there is a
primitive odd character $\xi_*\pmod{2^n}$.
Twisting a primitive character modulo $2^n$ by a character of
conductor $2^c$ with $c<n$ preserves primitivity.
For every primitive $\xi_*\pmod{2^n}$, odd $v>1$, odd positive $u$,
and $a\ge0$,
\[
 \xi_*(v)^{2^a u}=1\quad\Longrightarrow\quad
                     2^n<v^2 2^a.
 \tag{AX.10}
\]

Suppose $p$ is odd (allow $p=1$), $d>0$,
$4\mid\ell^d-1$, and $\omega_0,\theta_0,\xi_0$ have respective
levels $\ell,N,p$, with odd conductor for $\theta_0$ and orders
dividing $\ell^d-1$.  Let $\chi_2$ have conductor $2^c$, $c<n$.
If the product $\omega_0\chi_2\theta_0\xi_0\xi_*$ is odd and
\[
       p\ell\operatorname{cond}(\theta_0)
                  4^{v_2(\ell^d-1)}\le2^n,
 \tag{AX.11}
\]
then it equals $\operatorname{prodX}(\varphi,\xi_1)$ on units
of the common level, for literal $X\!\operatorname{Data}
(\ell,d,\varphi,\xi_1)$, with
$\operatorname{cond}(\varphi)\mid p\ell N$.
\end{proposition}

\begin{proof}
The standard generators give
$(\mathbb Z/2^n)^\times=\langle-1\rangle\times\langle5\rangle$,
where $5$ has order $2^{n-2}$.  To check the order, repeated squaring
shows $v_2(5^{2^j}-1)=j+2$; the subgroup $1+4\mathbb Z/2^n$
has exactly $2^{n-2}$ elements.  Send $-1$ to $-1$ and $5$ to
a primitive $2^{n-2}$-th root in $O$.  The resulting character is
odd and is nontrivial on $1+2^{n-1}$, the kernel of reduction to
level $2^{n-1}$, so it is primitive.  Conversely every primitive
character must be nontrivial on this order-two kernel, and takes
the value $-1$ there.  A lower-conductor twist is trivial on that
kernel; this proves the twisting assertion.

For (AX.10), the assertion for $n\le1$ is immediate from $v\ge3$.
For $n\ge2$, put $b=v_2(v^2-1)\ge3$.  Odd powering preserves
$v_2(x-1)$ for $x\equiv1\pmod4$, and each subsequent squaring
increases it by one.  Indeed
$(1+2^b t)^2-1=2^{b+1}(t+2^{b-1}t^2)$ with the parenthesis odd;
the odd-power statement follows from the geometric sum modulo $2$.
Thus
$v_2(v^{2u2^a}-1)=b+a$.
If $b+a<n$, further squaring produces a power of $v^{2^a u}$
congruent to $1+2^{n-1}$ modulo $2^n$.  Its character value is
$1$ under the assumed equality, contradicting primitivity.
Therefore $n\le b+a$, and
$2^n\le2^{b+a}<v^2 2^a$, as $2^b\le v^2-1$.

Take the product of $\omega_0,\theta_0,\xi_0$ at a common level
and pass to its primitive inducing character $\varphi$.  Its
conductor $f$ divides $p\ell\operatorname{cond}(\theta_0)$,
is odd, and its order divides $\ell^d-1$.  These assertions follow
directly by quotienting the finite unit group by the character
kernel: removing a redundant level preserves its image and order.
Inflate $\chi_2$ to $2^n$ and multiply by $\xi_*$ to obtain the
primitive $\xi_1$.  The conductors, the original parity equality,
and $f4^{v_2(\ell^d-1)}\le2^n$ now give every clause of the
$X\!\operatorname{Data}$ record in BR7.  Evaluation on a Galois
element is compatible with this identity: choose its one
cyclotomic residue modulo $p\ell N2^n$, represented by an integer
coprime to that level, and reduce the same integer at every factor.
This also proves the source \texttt{dirX} and \texttt{dirVal}
compatibility at the original $\Omega$.
\end{proof}

\subsection{Choices on the original algebraic closure}

\begin{proposition}[AX3: Frobenius choices and the prime-to-$\ell$ limit]
\label{prop:ax-choices}
For each rational prime choose a place $A_p$ of the original
$\Omega$ and an arithmetic Frobenius $\sigma_p$.  They may be
chosen so that $c(\sigma_\ell)=1$.  For any finite character
$\psi:G\to O^\times$ with open kernel there is
$s_0\in I_{A_\ell}\cap\ker\psi\cap\ker\omega$
with $c(s_0)$ of infinite order.

If $G_0$ is any profinite group, $g\in G_0$ and $\ell$ is prime,
the sequence $g^{\ell^{n!}}$ converges to an element $g^{(\ell')}$.
For any adically continuous lift $\rho$ of a finite-residue
representation, of any finite matrix size, suppose
$\bar\rho(g)^m=1$ with $(m,\ell)=1$.  Then
\[
 \bar\rho(g^{(\ell')})=\bar\rho(g),
 \qquad \rho(g^{(\ell')})^m=1
 \tag{AX.12}
\]
in its source complete Noetherian local coefficient algebra.
The residual continuity needed here follows from the lift.
Finally, any two nontrivial group
characters to groups have an element at which both are nontrivial.
These choices supply the literal \texttt{rawEx} output whenever
the two required residual characters are nontrivial.
\end{proposition}

\begin{proof}
Finite separable local extensions identify the chosen decomposition
group with the Galois group of $\mathbb Q_p$; the quotient by inertia
is $\widehat{\mathbb Z}$ with arithmetic Frobenius generator.  Choose
a lift.  At $\ell$, inertia maps onto
$\operatorname{Gal}(\mathbb Q_\ell(\mu_{\ell^r})/\mathbb Q_\ell)
\simeq(\mathbb Z/\ell^r)^\times$ for every $r$: the polynomial
$\Phi_{\ell^r}(1+X)$ is Eisenstein, so the finite cyclotomic
extension is totally ramified and has exactly that degree.
The nested closed fibers in the compact inertia group show that
$c(I_{A_\ell})=\mathbb Z_\ell^\times$, with the full inverse-limit
target, not only each separate quotient.  Multiply the chosen
Frobenius by an inertia element of cyclotomic value
$c(\sigma_\ell)^{-1}$.  This changes no residue Frobenius and gives
the asserted cyclotomic value.  At $p\ne\ell$, roots of unity of
order prime to $p$ lift from unramified residue extensions by
Hensel's lemma, so inertia fixes them and arithmetic Frobenius
acts by $\zeta\mapsto\zeta^p$.

Choose $u\in I_{A_\ell}$ with $c(u)=1+\ell$.
The open kernel of $\psi$ has finite index $b$; its image is
therefore finite of order $b$.  Then
$s_0=u^{b(\ell-1)}$ kills both $\psi$ and $\omega$.
The rational integer $1+\ell$ has infinite order in
$\mathbb Z_\ell^\times$ (injectivity of $\mathbb Z\to\mathbb Z_\ell$
would otherwise give $(1+\ell)^a=1$ in $\mathbb Z$ for some $a>0$).
Its indicated positive power still has infinite order.
Teichm\"uller lifting is the unique lifting of every element of
$\mathbb F_\ell^\times$ to a root of $X^{\ell-1}-1$ in
$\mathbb Z_\ell$; uniqueness proves multiplicativity.  This gives
$\omega^{\ell-1}=1$, $\bar\omega=\bar c$, and its factorization
through the cyclotomic action modulo $\ell$.
Complex conjugation extends from the algebraic real closure to
$\Omega$ and has cyclotomic value $-1$; all these choices are
made in the initially specified $G$.

In a finite quotient, write the order of $g$ as $\ell^a b$ with
$(b,\ell)=1$.  For all sufficiently large $n$, $\ell^{n!}$ is
$0$ modulo $\ell^a$ and $1$ modulo $b$, by Euler's theorem and
$\varphi(b)\mid n!$.  Thus the image of the sequence is eventually
constant in every finite quotient.  The constants are compatible,
so give an element of the profinite inverse limit and prove
convergence.  This proves the existence assertion without a
first-countability assumption on the group.

To prove the lift assertion let $\mathfrak m$ be the maximal ideal
of its coefficient algebra and $X=\rho(g)^m$.  Then
$X\equiv1\pmod{\mathfrak m}$ and $\ell\in\mathfrak m$.
The identity
\[
 X^\ell-1=(X-1)\left(\ell\,1+
                   \sum_{j=0}^{\ell-1}(X^j-1)\right)
 \tag{AX.13}
\]
shows that entrywise divisibility $X-1\in M_\nu(\mathfrak m^a)$,
$a\ge1$, implies $X^\ell-1\in M_\nu(\mathfrak m^{a+1})$
for any finite matrix size $\nu$.
Hence $X^{\ell^j}\equiv1\pmod{\mathfrak m^{j+1}}$.
Continuity modulo each $\mathfrak m^b$ and the convergence just
proved give $\rho(g^{(\ell')})^m\equiv1$ modulo every such power.
The ring is separated, proving (AX.12).  The residual identity
uses $\ell^{n!}\equiv1\pmod m$ and continuity of the finite
residual representation.

For the final assertion choose $x$ nontrivial under the first
character and $y$ nontrivial under the second.  If neither already
works, the first character is $1$ at $y$ and the second is $1$
at $x$, so $xy$ works.  Applied to $\bar\alpha,\bar\Psi$, this
chooses the same $g_0$ for both inequalities in (AX.8).
\end{proof}

\begin{proposition}[AX4: the residual local conditions and the generic good datum]
\label{prop:ax-generic}
Assume $e\ge2$, $\psi\in h.X(O,e)$ and $\psi_P\ne1$.
Let $p$ be odd, allowing $p=1$, and choose an even character
$\xi_0\pmod p$ and a primitive odd character $\xi_*\pmod{2^n}$.
Suppose
\[
\begin{gathered}
 v\nmid p\quad(v\in T_b),\qquad
 \xi_0^{\ell^{f_0}-1}=1,\qquad c_2<n,\\
 p\ell\operatorname{cond}(\theta)4^{n_0}\le2^n,
 \qquad u\text{ odd},\quad
 \psi^{2^a u}=\xi_0^{2^a u}=1,\qquad
 v^2 2^a\le2^n\quad(v\in T_b).
\end{gathered}\tag{AX.14}
\]
Then, using BR7, there is an actual $D$ on the original $\Omega,O$
with $D.\mathtt{Good}$ and
\[
 D.e=e,\quad D.\psi=\operatorname{dirG}\psi,\quad
 D.S=\operatorname{primeFactors}(2\ell Np),\quad
 D.\xi=\operatorname{dirG}\xi_0\operatorname{dirG}\xi_*.
 \tag{AX.15}
\]
It has the literal $\mathtt{OfChar}(e,\psi,n,p,
\operatorname{ord}\xi_0)$ witnesses.
\end{proposition}

\begin{proof}
Take the character data and the original-field choices of AX3.
For now postpone only $g_0$; AX1 already gives $T\subset T_b$.
Reduction kills $\kappa$, since $k^\times$ has order prime to
$\ell$.  With $\omega$ also read as the Dirichlet character
modulo $\ell$, one has
\[
 \bar\alpha=\overline{\xi_0\xi_*},\qquad
 \bar\Psi=\overline{\omega^e\psi_2\theta}.
 \tag{AX.16}
\]
For every $v\in T$, $\bar\omega(\sigma_v)=1$.
At $\ell$ this is AX3; at every other $v$ it follows from
$v\equiv1\pmod\ell$ and the Frobenius rule.
If $\overline{\xi\Psi^i}(\sigma_v)=1$, raise the equality to
$2^a u$.  The factors $\psi$ and $\xi_0$ disappear by (AX.14),
and the remaining cyclotomic residue is $1$.  Thus
$\overline{\xi_*(v)^{2^a u}}=1$.
Reduction is injective on roots of unity of order prime to $\ell$:
if $z\equiv1$ and $z^m=1$, $(m,\ell)=1$, then
$1+z+\cdots+z^{m-1}$ is a unit, so
$(z-1)(1+\cdots+z^{m-1})=0$ forces $z=1$.
The dyadic root in question therefore equals $1$ in $O$.
AX2 gives $2^n<v^2 2^a$, a contradiction.  This proves the full
integer-parameter condition (AX.7), including negative $i$.

We next verify the six residual vanishing conditions, with their
actual local subgroups.  The prime-to-$\ell$ residual diagonal
character group has order dividing $(|k|-1)^4$; consequently
$K=\ker\bar r_0$ is normal of finite index prime to $\ell$.
Let $z$ satisfy the defining conditions of $X\!\operatorname{Van}
(\gamma)$, where $\bar\gamma$ is a character covered by BR7.
At a place away from $T$, choose $H_A=K$ in BR7's
$\operatorname{UnrEv}$.  Its relative index in
$I_A\cap\ker\bar\gamma$ divides $[G:K]$, since $K$ is normal,
and the required vanishing is exactly $\operatorname{UnrT}$.
At a place above $v\in T$, conjugate inertia into the selected
$I_{A_v}$.  The local coboundary is zero on the kernel of
$\bar\gamma$.  For such an inertia element $x$, the cocycle
identity gives, explicitly,
\[
 z(g^{-1}xg)=\bar\gamma(g)^{-1}z(x)
       \quad\text{if }\bar\gamma(x)=1.
 \tag{AX.17}
\]
Thus $z$ vanishes on $I_A\cap\ker\bar\gamma$ and one may take
$H_A=G$.  This verifies the original relative-index condition at
every valuation subring.  BR7 gives the global coboundary in the
original field $k$.  Applying the same argument to the inverse
character gives its inverse case as well.

To apply it, use AX2 with the following three rows:
\[
\begin{array}{c|ccccc}
\gamma&\omega_0&\chi_2&\theta_0&\text{prime character}&
                 \text{dyadic character}\\ \hline
\alpha&1&1&1&\xi_0&\xi_*\\
\beta&\omega^e&\psi_2&\theta&\xi_0^{-1}&\xi_*^{-1}\\
\vartheta&\omega^e&\psi_2&\theta&\xi_0&\xi_*
\end{array}\tag{AX.18}
\]
Take $d=f_0$ in all rows.  AX1 and (AX.14) give the order and
conductor inequalities; for the first row use
$1\le\operatorname{cond}\theta$.  The parity of
$\omega^e\psi_2\theta$ is even, because
$\kappa(-1)=1$ and $\psi(-1)=(-1)^e$.
Each row is therefore odd.  Formula (AX.16), its inverse, and
the common-residue comparison in AX2 identify the residual
characters exactly.  All six applications use the same $D$, $K$
and originally chosen places; no new residual coefficient field
has appeared.

For $g_0$, (AX.7) with $v=\ell$ and $i=0$ proves
$\bar\alpha\ne1$.  Also $\bar\Psi\ne1$.  Here is the needed
argument rather than an assumption of that fact.  The prime-to-$\ell$
part of $\psi$ is the product of $\psi_P$, unramified at $\ell$,
and an $\ell$-local character.  The latter factors through
$(\mathbb Z/\ell)^\times$, since its restriction to the
principal-unit $\ell$-group is trivial.  It is some power
$\omega^b$.  If $\bar\Psi=1$, injectivity of reduction on
prime-to-$\ell$ roots gives $\omega^{e+b}\psi_P=1$ in $O$.
Restriction to $I_{A_\ell}$ and cyclotomic surjectivity force
$\omega^{e+b}=1$.  Hence $\psi_P=1$, contrary to hypothesis.
The finite cyclotomic action on the original $G$ is surjective,
so this last conclusion is also an equality of the original
Dirichlet characters.  AX3 now provides $g_0$.

It remains to check the other fields of \texttt{Good}.
The set in (AX.15) contains $2,\ell$ and only primes; its data
have the stated $e$, selected places and complex conjugation.
Finite Dirichlet characters have open kernel and are unramified
away from the primes in their levels, by AX3.
The order of $\xi_0\xi_*$ divides
$(\ell^{f_0}-1)\varphi(2^n)$, which is prime to $\ell$.
No prime of $T$ divides $p2^n$, so $\xi$ is also unramified
there.  Its parity is odd, whereas $\Psi(\mathrm{cc})=1$.
The Teichm\"uller, Frobenius, limit and $s_0$ clauses are exactly
AX3.  Together with (AX.7), (AX.17)--(AX.18) and the chosen $g_0$,
these are all the fields listed above.  Finally $\xi_*,\xi_0$
themselves witness \texttt{OfChar}, with precisely their original
levels and $\operatorname{ord}\xi_0$.
\end{proof}

\begin{proposition}[AX5: literal dataEx in both original host cases]
\label{prop:ax-dataex}
Let $h.\mathtt{Big}(t)$, $0<e<t$, $\ell\ne2$ and
$\psi\in h.X(O,e)$.  Suppose Case II or Case III holds and the
original $O$ contains all roots of unity of orders
$2^{n_{\rm II}}$ and $2^{n_{\rm III}}$.  Then there is
$D:\mathtt{Raw}(\Omega,O)$ with $D.\mathtt{Good}$ and the literal
$\mathtt{OfCase}(D,h,e,\psi)$.  More precisely, the parameters
$(n',p,m_0)$ in \texttt{OfChar} may be taken as
\[
 (n_{\rm II},1,1)\quad\text{in Case II},\qquad
 (n_{\rm III},p_{\!1},\operatorname{ord}\psi)
                         \quad\text{in Case III}.
 \tag{AX.19}
\]
\end{proposition}

\begin{proof}
In Case II choose the primitive odd dyadic character at the
prescribed $n=n_{\rm II}$ by AX2, set $p=1$, $\xi_0=1$, and write
$\kappa^{\ell^b}=1$.  Use
\[
 a=c_2+1+v_2(o),\qquad u=\operatorname{oddpart}(o)\ell^b.
 \tag{AX.20}
\]
The integer $u$ is odd and
$2^a u=2^{c_2}(2o)\ell^b$ kills every factor of (AX.2).
It also kills $\xi_0$.  The first maximum in (AX.3) gives
$\ell\operatorname{cond}(\theta)4^{n_0}\le2^n$.
For $v\in T_b$, AX1 gives $v\le\ell q$; hence
\[
 v^2 2^a\le\ell^2q^2\,2^{c_2+1}o
            \le2^{c_2+6}o\ell^3q^2\le2^n.
 \tag{AX.21}
\]
All other conditions of (AX.14) follow from AX1 or $p=1$.
AX4 supplies the first parameter choice in (AX.19).

In Case III, (AX.9) applies.  Put $p=p_{\!1}$ and
$m=\operatorname{ord}\psi$.  Since $m\mid\operatorname{oddpart}(o)$
and $p\equiv1\pmod{\operatorname{oddpart}(o)}$, choose a character
$\xi_0$ of $(\mathbb Z/p)^\times$ of order exactly $m$.
There is already a root of order $m$ in $O$: the finite image of
$\psi$ in the multiplicative group of a field is cyclic of order
$m$.  Send a generator of $(\mathbb Z/p)^\times$ to that root.
Thus no extra roots are required.  The order $m$ is odd, so
$\xi_0(-1)=1$.  Choose the primitive odd dyadic character at
$n=n_{\rm III}$; set $a=0,u=m$.  Then $T_b=\{\ell\}$,
$\ell\ne p$, and
\[
 p\ell\operatorname{cond}(\theta)4^{n_0}
 \le p\ell q4^{n_0}\le2^n,
 \qquad \ell^2\le8\ell^2\le2^n.
 \tag{AX.22}
\]
The order of $\xi_0$ divides $\ell^{f_0}-1$, because
$\theta=\psi$.  All conditions of AX4 now hold, giving the second
case with $m_0=m$.  This is the full quantified \texttt{dataEx}
implication with BR7 installed, on the original $\Omega,O$.
\end{proof}

\subsection{All local facts for any good datum}

\begin{proposition}[AX6: literal facts and the normalization unit]
\label{prop:ax-facts}
For every actual $D:\mathtt{Raw}(\Omega,O)$, $D.\mathtt{Good}$
implies all clauses of \texttt{facts} in \texttt{LMemMax}:
\begin{enumerate}
\item If $p\in S_b$, then $p$ is prime, $p\ne\ell$,
$\ell\mid p-1$, and $\omega|_{G_{A_p}}=1$.
\item For $v\in T$, $\sigma_v\in G_{A_v}$,
$\bar\omega(\sigma_v)=1$, and
$\bar p_i(\sigma_v)\ne\bar p_j(\sigma_v)$ whenever
$\lambda_i\ne\lambda_j$, where $\lambda=(1,0,1,0)$.
\item $\operatorname{UnrT}(p_i,1)$ holds for all $i$.  On
$I_{A_\ell}\cap\ker\psi\cap\ker\omega$ one has $p_i=c^{w_i}$.
For $p\in S_b$ and $u\in I_{A_p}$,
$\alpha(u)=1$ and $\Psi(u)=\vartheta(u)=\psi(u)$.
\item $p_i p_{3-i}=\vartheta$, the four values $p_i(s_0)$ are
distinct, and $\alpha(\mathrm{cc})=\beta(\mathrm{cc})=
\vartheta(\mathrm{cc})=-1$.
\item $\phi\in G_{A_\ell}$ and
$\bar r_0(\phi)=\bar r_0(\sigma_\ell)$.  Each $p_i$ is adically
continuous; $\ker\bar r_0$ and $\ker\omega$ are open.
Each $p_i$ is unramified outside $S$.
\end{enumerate}
In addition the characters $\alpha,\beta,\vartheta$ on the common
inertia subgroup in item 3 have positive weights $1,e-1,e+1$.
For every actual complete-local lift $\rho$ of $\bar r_0$, put
\[
 H(\rho)=\rho(g_0)+\vartheta(g_0)\rho(g_0^{-1}).
\]
Then $H(\rho)_{00}-H(\rho)_{22}$ is a unit, with exact residue
\[
 \overline{H(\rho)_{00}-H(\rho)_{22}}
       =(\bar\alpha(g_0)-1)(\bar\Psi(g_0)-1).
 \tag{AX.23}
\]
Every algebraic specialization of this unit is nonzero.
\end{proposition}

\begin{proof}
This proposition uses only \texttt{Good}, so also applies when
$D$ was not constructed by AX5.  First $\omega$ is the unique
Teichm\"uller character: its values and those of the character
in AX3 are prime-to-$\ell$ roots with identical reductions.
The injectivity argument in AX4 proves their equality.  In
particular it is unramified away from $\ell$ and has open kernel.

The open kernel of $\psi$ gives a finite abelian image.  In a
finite abelian extension of $\mathbb Q_p$, the prime-to-$p$ part
of inertia has exponent dividing $p-1$.  Here is the precise local
fact needed.  Wild inertia is pro-$p$.  In the tame quotient inertia
is generated topologically by compatible roots of a uniformizer;
arithmetic Frobenius conjugates a tame generator $\tau$ to
$\tau^p$.  This identity follows by lifting the action
$\zeta\mapsto\zeta^p$ on the prime-to-$p$ roots of unity used in
$\tau(\pi^{1/m})=\zeta_m\pi^{1/m}$.  After passing to an abelian
finite quotient it gives $\tau^{p-1}=1$.  Thus an $\ell$-divisor
of its inertia order with $p\ne\ell$ divides $p-1$.
Applying this to the actual image of $\psi|_{G_{A_p}}$ proves
$\ell\mid p-1$ for $p\in S_b$.
The Teichm\"uller character is trivial on $I_{A_p}$ and takes
the value $1$ at Frobenius because $p\equiv1\pmod\ell$.
Inertia and a Frobenius lift generate a dense subgroup of the
decomposition group, so $\omega=1$ on all of $G_{A_p}$.

The second assertion follows from the Frobenius clause of
\texttt{Good}, the preceding paragraph, and $c(\sigma_\ell)=1$.
Put $a=\bar\xi(\sigma_v)$ and $b=\bar\Psi(\sigma_v)$.
The four residual values are $(ab,b,a,1)$; condition (AX.7) says
$ab^j\ne1$ for every integer $j$.  The four cross-$\lambda$
comparisons are $ab\ne b$, $ab\ne1$, $a\ne b$ and $a\ne1$.
They follow respectively from $j=0,1,-1,0$.  This lists every
comparison, including their reverses.

For $\operatorname{UnrT}$, a place away from $T$ is away from
$\ell$.  There $c=\omega=1$ on inertia, so
$\alpha=\xi$ and $\Psi=\psi$.  The residual-kernel condition
gives $\bar\xi=\bar\psi=1$.  The order of $\xi$ is prime to
$\ell$, so $\xi=1$ by AX4.  If the place is away from $S$,
$\psi=1$ directly.  Otherwise it lies above some $p\in S$
with $p\notin S_b$; conjugacy of places identifies the inertia
image with the selected one, whose order is prime to $\ell$.
Thus $\psi=1$ there as well.  All $p_i$ equal $1$.
On the common inertia subgroup above $\ell$, both $\psi$ and
$\omega$ are $1$, and $\xi=1$ by \texttt{Good}; substitution in
(AX.5) gives the weights, including the three positive ones.
For $p\in S_b$, $\xi$ is unramified and the preceding
$c=\omega=1$ identities give item 3.

The pairing identities $p_i p_{3-i}=\vartheta$ follow from
(AX.5).  Distinctness at $s_0$ follows from the infinite order
of $c(s_0)$ and the four distinct weights for $e\ge2$.
Complex conjugation has $c=\omega=-1$, while $\xi=-1$ and
$\Psi=1$, proving the three oddness assertions.
The closed decomposition group contains all powers of
$\sigma_\ell$, hence their limit $\phi$.
Reduction modulo any power of $\mathfrak m_O$ of a cyclotomic
character is locally constant: a sufficiently deep cyclotomic
principal subgroup maps into that ideal power.  The finite
characters have open kernels.  Products and powers therefore
give the asserted adic continuity of all $p_i$, and reduction
of the diagonal representation has open kernel.
Its image at $\sigma_\ell$ has order dividing $|k|-1$.
AX3 proves $\bar r_0(\phi)=\bar r_0(\sigma_\ell)$.
Outside $S$, the finite characters and the cyclotomic character
are all unramified, since $\ell\in S$.  This finishes every
clause of the literal source \texttt{facts}.

Finally write $a=\bar\alpha(g_0)$ and $b=\bar\Psi(g_0)$.
The diagonal entries $0,2$ of the residual $H$ are respectively
$ab+1$ and $a+b$.  Their difference is $(a-1)(b-1)$, which is
nonzero by \texttt{Good}.  An element of a local ring with
nonzero residue is a unit.  Ring maps carry units to units,
so its image in $\overline{\mathbb Q}_\ell$ is nonzero.  This
proves (AX.23), with no extra normalization premise on $\rho$.
\end{proof}

\begin{proposition}[AX7: the literal common level and exponent]
\label{prop:ax-level-integer}
Suppose the same good $D$ has
$\mathtt{OfChar}(e,\psi,n',p,m_0)$, with $c_2\le n'$ and $m_0\mid o$.
Define
\[
             L=\ell 2^{n'}Np,
             \qquad k_0=(\ell-1)2^{n'}(2o).
             \tag{AX.24}
\]
Then $L,k_0>0$, each of $D.\psi,D.\xi,D.\omega$ has $k_0$-th
power one, and each satisfies the literal
$\operatorname{KillsMu}(L)$: it is $1$ on every $g$ fixing all
$L$-th roots of unity in $\Omega$.  For the data of AX5,
$D.S=\operatorname{primeFactors}(2\ell Np)$ has exactly the
prime support required by G1, and the parameters are (AX.19).
\end{proposition}

\begin{proof}
The levels $N,p,2^{n'}$ divide $L$.  Their cyclotomic actions
are therefore trivial whenever the action on $\mu_L$ is trivial.
The Teichm\"uller character factors through $\mu_\ell$, also
contained in $\mu_L$.  This proves all three \texttt{KillsMu}
assertions on the same $G$.

The dyadic component of $\psi$ has order dividing $2^{c_2}$,
and every odd component has order dividing $2o$, by the CRT
argument of AX1.  Hence $\psi^{2^{n'}(2o)}=1$.
The prime factor $\xi_0$ has order $m_0\mid o$ and the dyadic
factor has order dividing $\varphi(2^{n'})\mid2^{n'}$.
Thus $\xi^{2^{n'}o}=1$.  Finally $\omega^{\ell-1}=1$.
Each exponent divides $k_0$, proving the literal \texttt{levOf}
claims.  Positivity follows from $\ell>2$, $o,N,p>0$.
AX5 supplies $m_0=1$ in Case II and
$m_0=\operatorname{ord}\psi\mid\operatorname{oddpart}(o)$ in
Case III, together with the required $c_2\le n'$.
The assertion about $S$ is the exact equality (AX.15).
\end{proof}

\subsection{Ordinary flags on the component containing the diagonal point}

The source's local condition is a polynomial condition, whose
meaning matters under specialization.  For scalars $a_i$ put
\[
 P_i(X)=\prod_{j\ne i}(X-a_j1),\qquad
                  d_i=\prod_{j\ne i}(a_i-a_j).
 \tag{AX.25}
\]
The condition $\operatorname{Ord}(I,(p_i),s_0,r)$ says that
$r(s_0)$ has characteristic polynomial $\prod_i(T-p_i(s_0))$,
that $P_j(r(s_0))r(g)P_i(r(s_0))=0$ for $i<j$, and that
$P_i(r(s_0))(r(u)-p_i(u)1)P_i(r(s_0))=0$ for $u\in I$.
These are exactly the three source fields, with unnormalized
projectors, so they make sense even before $d_i$ is invertible.

An actual point of the ordinary component is a specialization
of a lift $\rho:G_{A_\ell}\to M_4(B)$, where $B$ is a source
complete Noetherian local $O$-algebra which is a domain,
$\rho$ satisfies the similitude and ordinary equations, and there
is an $O$-algebra map $x:B\to O$ carrying $\rho$ to
$r_0|_{G_{A_\ell}}$.  A second map
$y:B\to\overline{\mathbb Q}_\ell$ gives the point.
An intervening coefficient algebra in \texttt{PtOn}/\texttt{OnComp}
is absorbed by composing its map with $y$; the domain $B$ and
the base point $x$ are retained.

\begin{proposition}[AX8: the full ordinary specialization theorem]
\label{prop:ax-ordinary}
Let $D$ be good.  For every actual
$r:G_{A_\ell}\to M_4(\overline{\mathbb Q}_\ell)$ on the source
ordinary component, there is a basis making every $r(g)$ upper
triangular with diagonal continuous characters $\kappa_i$, such
that
\[
\begin{gathered}
 \kappa_i(u)=p_i(u)\quad(u\in I_{A_\ell}),\qquad
 |\kappa_i(g)-p_i(g)|_\ell<1\quad(g\in G_{A_\ell}),\\
 \kappa_0\kappa_3=\kappa_1\kappa_2=\vartheta.
\end{gathered}\tag{AX.26}
\]
For every open subgroup $U\le G_{A_\ell}$ containing
$\sigma_\ell$ and every $i\ne j$, there is $g\in U$ with
$\kappa_i(g)\ne c(g)\kappa_j(g)$.  The same nonidentity holds
on $G_{A_\ell}$ also for $i=j$.  These assertions supply the
literal \texttt{spPtL} interface, including its uniformity over
all such open subgroups.
\end{proposition}

\begin{proof}
We first prove the specialization estimate used in the argument.
Let $R$ be a local $O$-algebra with $\mathfrak m_O R\subset
\mathfrak m_R$, and let $y:R\to\overline{\mathbb Q}_\ell$ be an
$O$-algebra map.  Then
\[
 |y(a)|_\ell\le1\ (a\in R),\qquad
 |y(a)|_\ell<1\ (a\in\mathfrak m_R).
 \tag{AX.27}
\]
For the first assertion, if $|y(a)|>1$, the inverse of $y(a)$
has a monic minimal polynomial over $\operatorname{Frac}(O)$
whose nonleading coefficients lie in $\mathfrak m_O$.
Indeed all conjugates have the same norm less than $1$, by
uniqueness of the extension of the complete-field valuation;
their elementary symmetric functions lie in $O$ and have norm
less than $1$.  Multiplying the polynomial identity by a power
of $y(a)$ gives $1+b_1y(a)+\cdots+b_ny(a)^n=0$ with
$b_i\in\mathfrak m_O$.  The corresponding element
$1+b_1a+\cdots+b_na^n$ of $R$ is a unit and cannot map to zero.
If $a\in\mathfrak m_R$ and $|y(a)|=1$, instead use its monic
minimal polynomial.  All coefficients are in $O$ and the constant
coefficient is a unit.  Evaluating the polynomial at $a$ again
gives a unit of $R$ mapping to zero.  This proves (AX.27).
The arguments apply also to zero images, which already satisfy
the asserted inequalities.

If $R$ is Noetherian, choose generators $a_1,\ldots,a_b$ of
$\mathfrak m_R$ and put $a=\max_i|y(a_i)|<1$; the zero ideal
case is immediate.  Products of generators and (AX.27) give
$|y(z)|\le a^n$ for $z\in\mathfrak m_R^n$.  Thus an adically
continuous entry function becomes an ordinary continuous
$\overline{\mathbb Q}_\ell$-valued function.  In particular all
specialized matrix entries below are continuous.

We also need the following precise divided-eigenvalue argument.
Let $B$ be the domain above, $d\in O\setminus\{0\}$,
$t\in B$, $x(t)=0$, and let $v=t/d\in\operatorname{Frac}(B)$
be integral over $B$.  The base point $x$ is a retraction on
$O$, so $O\to B$ is injective and this denominator is nonzero.
The algebra $B[v]$ is finite over the
Henselian local domain $B$.  It is a domain and hence a local
ring by the stated finite-algebra theorem.  The maps $x,y$
extend to the localization $B[1/d]$, with targets
$\operatorname{Frac}(O)$ and $\overline{\mathbb Q}_\ell$.
Restrict them to $B[v]\subset B[1/d]$.  The first sends $v$ to
zero, so $v$ is not a unit and lies in the maximal ideal of
$B[v]$.  Its maximal ideal contracts to that of $B$, since the
extension is finite integral.  Applying (AX.27) to the second
map proves
\[
                      |y(t)/d|_\ell<1.
 \tag{AX.28}
\]
This is the point where the same domain component and its actual
diagonal base point are used; a bare ordinary specialization
would not justify the conclusion.

Now set $a_i=p_i(s_0)$.  AX6 makes them distinct.  Over either
$\operatorname{Frac}(B)$ or $\overline{\mathbb Q}_\ell$ the
operators $d_i^{-1}P_i(r(s_0))$ are pairwise orthogonal rank-one
idempotents summing to $1$.  This follows by evaluating the
Lagrange interpolation identities modulo
$\prod_i(T-a_i)$ and using the characteristic polynomial; each
eigenspace is one-dimensional because its root is simple.
Choose one vector from each eigenline, in index order.
The stable equations imply that the sums of the first $i+1$
lines are invariant under every $r(g)$.  The matrices are
therefore upper triangular.  Their diagonal entries are
characters $\kappa_i$, since multiplication of upper triangular
matrices multiplies the diagonal entries, and their inverse
values exist from $r(g^{-1})$.  The inertia equations give
$\kappa_i(u)=p_i(u)$ exactly.  Moreover
\[
 d_i\kappa_i(g)=\operatorname{tr}\big(P_i(r(s_0))r(g)\big).
 \tag{AX.29}
\]

Use the same argument first over $\operatorname{Frac}(B)$.
Its diagonal value $v_i(g)$ is integral over $B$, since it is
a root of the monic characteristic polynomial of $\rho(g)$.
Write $t_i(g)=\operatorname{tr}(P_i(\rho(s_0))\rho(g))$.
Then $d_i v_i(g)=t_i(g)$ and at the actual diagonal point
$x(t_i(g))=d_i p_i(g)$.
Apply (AX.28) to the integral element $v_i(g)-p_i(g)$,
with numerator $t_i(g)-d_i p_i(g)$.  Equation (AX.29) after
specialization identifies its image with
$\kappa_i(g)-p_i(g)$, and proves the strict estimate in (AX.26).
Continuity of $\kappa_i$ follows from continuity of the entries
of $r$ after a fixed conjugation, or directly from (AX.29).

The original alternating form is preserved up to $\vartheta$.
In the chosen eigenbasis for $r(s_0)$ its matrix can have a
nonzero entry at $(i,j)$ only if
$a_i a_j=\vartheta(s_0)$.
Since $a_i=c(s_0)^{w_i}$ with $c(s_0)$ of infinite order,
this means $w_i+w_j=e+1$.
For the four weights and $e\ge2$ this occurs precisely when
$j=3-i$.  Nondegeneracy makes all four antidiagonal entries
nonzero.  For an upper triangular $r(g)$, the $(i,3-i)$ entry
of $r(g)^tJr(g)=\vartheta(g)J$ then gives
$\kappa_i(g)\kappa_{3-i}(g)=\vartheta(g)$.
This proves the exact two pairing identities in (AX.26).

Finally suppose $\kappa_i=c\kappa_j$ on an open subgroup $U$
containing $\sigma_\ell$.  The subgroup has finite index in
the compact decomposition group, so $s_0^b\in U$ for some
positive $b$.  Evaluation there, using inertia and infinite
order, gives $w_i=w_j+1$.  This is impossible for $i=j$.
For distinct indices the only pairs allowed by the four weights
are $(0,1),(2,3)$ and, when $e=2$, $(1,2)$.
Every such pair has $\lambda_i\ne\lambda_j$.
Evaluation at $\sigma_\ell$ gives
$\kappa_i(\sigma_\ell)=\kappa_j(\sigma_\ell)$ because
$c(\sigma_\ell)=1$.  The strict estimates in (AX.26) imply
$\bar p_i(\sigma_\ell)=\bar p_j(\sigma_\ell)$: the difference
of their values has norm less than $1$, which is precisely
membership in $\mathfrak m_O$.  This contradicts AX6.
Thus the required element of $U$ exists for every pair.
\end{proof}

\subsection{The full admissible class on the same data}

\begin{proposition}[AX9: literal level, dual twist and safety]
\label{prop:ax-admissible}
Let $D$ be any good datum and $M\ge1$.  Suppose
$\operatorname{KillsMu}(M)$ holds for $\psi,\xi,\omega$.
For every actual $D.\mathtt{EPt}$ representation $r$,
\[
 \mathtt{Level}\bigl(\ell,M,(e+1,e,1,0),r\bigr),
 \qquad \mathtt{IsCycWt}(\ell,e+1,\vartheta),
 \qquad \mathtt{DualTw}(r,\vartheta)
 \tag{AX.30}
\]
hold, with the original coefficient embedding into
$\overline{\mathbb Q}_\ell$.  Also $\ell\mid M$ and
$r$ satisfies $\mathtt{UnlAll}$, hence the first alternative of
the literal $\mathtt{SafeRA}$ condition.  This proves
\texttt{level}, \texttt{spUnl}, \texttt{teichDvd}, and
\texttt{safe4'} on their actual quantified inputs.
\end{proposition}

\begin{proof}
An actual point is the image under $R\to\overline{\mathbb Q}_\ell$
of an adically continuous lift over a source complete local
$O$-algebra.  Formula (AX.27) and the finite-generator estimate
in AX8 make its matrix entries continuous.  Since
$r(g)^{-1}=r(g^{-1})$, this is also continuity as a
general-linear-group representation.  Its similitude equations,
the equality $r=1$ on the prescribed inertia kernel outside $T$,
and its equations at every $p\in S_b$ survive specialization:
they are matrix polynomial equalities.  For the \texttt{OnComp}
local conditions, compose with the actual domain-component map.
This explains the needed part of the source \texttt{ptIn}
transport without assuming a representability theorem.

We prove the away-from-$\ell$ part of \texttt{Level} with global
characters, as its definition requires.  Let
$\varphi=(\psi,\xi,\omega):G\to A$ with $A$ its actual finite
abelian image.  Choose a place over a given rational prime
$p\ne\ell$.  If it is away from $T$, then for
$g\in I_A\cap\ker\varphi$ the cyclotomic character is $1$,
so $\bar r_0(g)=1$ and the point's unramified condition gives
$r(g)=1$.  If it lies above a member of $T$, that member is
some $p\in S_b$; choose the specified $A_p$ instead.
The point's actual $\mathtt{condP}$ equation gives
$r=1$ on $I_{A_p}\cap\ker\psi$, in particular on
$I_{A_p}\cap\ker\varphi$.

Thus in either case the inertia representation factors through
the finite abelian subgroup $\varphi(I_A)\le A$.
Every representing matrix has finite order and is diagonalizable
in characteristic zero, since its minimal polynomial divides a
separable polynomial $T^m-1$.  Commuting diagonalizable matrices
have a common eigenbasis: split into eigenspaces of one nonscalar
matrix and induct on their dimensions; if all are scalar the
claim is immediate.  The diagonal characters of $\varphi(I_A)$
extend to $A$.  Explicitly, when adjoining one generator $a$ to
a subgroup and its coset has order $b$, choose in the algebraic
closure a $b$-th root of the already specified value at $a^b$;
the rule on $ha^j$ is then well-defined and multiplicative.
Iterate through generators of the finite group.  Compose the
four extended characters with $\varphi$ to obtain four global
characters on the original $G$, diagonalizing inertia.  They
are all trivial on $\ker\varphi$, hence satisfy
$\operatorname{KillsMu}(M)$.  This is the complete
$\mathtt{LevelAway}$ assertion; it does not merely produce
local inertia characters.

At $\ell$, AX8 supplies the actual ordinary flag and characters
$\kappa_i$.  Define the following global finite characters:
\[
 (\chi_0,\chi_1,\chi_2,\chi_3)
       =(\omega^{-1}\xi\psi,\psi,\omega^{-1}\xi,1).
 \tag{AX.31}
\]
They all satisfy $\operatorname{KillsMu}(M)$, and
$p_i=c^{w_i}\chi_i$.  The inertia formula in (AX.26) is
therefore exactly the formula in $\mathtt{LevelAt}$.
The last part of AX8 gives
$\kappa_i\ne c\kappa_j$ for every pair, so we use the first
alternative of that definition.  No potentially crystalline
alternative is needed.  Over $\mathbb Q$ there is only one
finite place above $\ell$, and the chosen $A_\ell$ lies above
it; the comparison with the source prime ideal is just
$\mathcal O_{\mathbb Q}=\mathbb Z$ and its ideal $(\ell)$.
This proves $\mathtt{Level}$.

The multiplier is
$\vartheta=c^{e+1}(\omega^{-1}\xi\psi)$, whose finite part
satisfies $\operatorname{KillsMu}(M)$.  This is the actual
$\mathtt{IsCycWt}$ witness with the same $M$.
For the duality witness take the fixed matrix
\[
 J=\begin{pmatrix}0&0&0&1\\0&0&1&0\\0&-1&0&0\\-1&0&0&0\end{pmatrix},
 \qquad J^{-1}=-J.
 \tag{AX.32}
\]
The source adjoint is $J^{-1}X^tJ$; its similitude equation is
therefore exactly $r(g)^tJr(g)=\vartheta(g)J$.
This supplies the original $\mathtt{DualTw}$, with no change of
Frobenius or cyclotomic convention.

To see that $\ell\mid M$, suppose the contrary.  Inertia at
$A_\ell$ fixes $\mu_M$, so $\omega=1$ there by
$\operatorname{KillsMu}(M)$.  But AX3 supplies an inertia element
with $c=-1$, whose Teichm\"uller value is $-1$.  Since $\ell$
is odd this is a contradiction.  The noncritical ordinary flag
just proved holds for every place above $\ell$, not only the
selected one: if $A'=gA_\ell$, transport the flag by $r(g)$
and the characters by $u\mapsto g^{-1}ug$.
The global cyclotomic character is invariant under conjugation,
so all inequalities $\kappa_i\ne c\kappa_j$ are preserved.
The transport map on decomposition groups is an isomorphism,
so no nonidentity is lost.  This is precisely
$\mathtt{UnlAll}$.  Combining it with $\ell\mid M$ proves
the first clause of $\mathtt{SafeRA}$, and hence the literal
\texttt{safe4'} conclusion.
\end{proof}

\begin{proposition}[AX10: the complete G1 implication]
\label{prop:ax-g1}
Assume the image-density hypothesis, with the literal
\texttt{SymplecticFourImageDense} statement.
For every original host and original Case II or Case III
character in AX5, adjoining the roots already allowed by SP8
produces an actual $D$ satisfying every clause of G1 in SP10.
Its parameters and level are exactly (AX.19) and (AX.24).
Moreover $D.\mathtt{Fin0}$ and $D.\mathtt{NoEnd}$ hold, and there
is one embedding
$\iota:\overline{\mathbb Q}_\ell\simeq\mathbb C$, chosen before
the varying point, such that every actual irreducible
$D.\mathtt{EPt}$ representation belongs to
\[
 \mathtt{AdmRAlt}\bigl(\iota,L,(e+1,e,1,0),e+1\bigr).
 \tag{AX.33}
\]
There is no remaining construction, character, residue-field,
level, duality, ordinary-specialization or safety premise in
this implication.  The actual image-density supplier itself
remains an arithmetic obligation.
\end{proposition}

\begin{proof}
First suppose the original $O$ already has the two required
sets of dyadic roots.  AX5 constructs $D$ on its original
$\Omega,O$ and the exact \texttt{OfCase}.  AX1 and AX5 provide
$c_2\le n'$ and $m_0\mid o$, so AX7 gives the common $L,k_0$
and all three $\mathtt{KillsMu}$ and power assertions.
AX9 applies to every actual point of this same $D$, giving
its level, multiplier of weight $e+1$, dual twist and safety.
Now apply \texttt{SymplecticFourImageDense} to this $D$, with
the same host, $t,e,\psi$ and original \texttt{Good}/\texttt{OfCase}
witnesses.  Its conclusion is exactly
\[
 D.\mathtt{Fin0}\ \wedge\ D.\mathtt{NoEnd}\ \wedge\
 \exists\iota\ \forall r\in D.\mathtt{EPt},\quad
 \mathtt{Irr}(r)\Longrightarrow\mathtt{IsAutQA}(\iota,r).
 \tag{AX.34}
\]
Here $\mathtt{Irr}$ has the source meaning that the linear span
of the matrices $r(g)$ is all of $M_4$.  It implies the
$\mathtt{Af.IsIrr}$ field of the admissible record: an invariant
subspace would be invariant under every matrix, and a nonzero
vector can be carried to every vector by suitable matrices.
Thus the subspace is the whole vector space.  This is the exact
direction used for the actual points in SP10.
The quantifiers in (AX.34), rather than separately chosen
embeddings for each point, supply the common $\iota$.
Irreducibility, the stated automorphy assertion, and the four
clauses already proved are precisely the fields of
$\mathtt{AdmRAlt}$ in (AX.33).

For a coefficient DVR initially lacking the roots, use the
finite extension allowed in SP8 and RV1.  Its integral closure
inside the same $\overline{\mathbb Q}_\ell$ is a finite
coefficient DVR containing the two finite sets of roots.
The injective coefficient map preserves the original
character's kernel, order, conductor, CRT components, parity,
$f_0,n_0,c_2$, and hence both integers in (AX.3).
These facts can also be read directly: equality of two
character values is unchanged by an injective ring map,
so every finite quotient through which it factors is unchanged.
Construct $D$ over this allowed extension and perform the
preceding argument there.  This does not assert that the
new $D$ is a datum over the smaller coefficient ring.
The final normalized Selmer-length comparison is returned to
the original coefficient ring by the exact ramification-factor
cancellation already proved in RV1 and SP8.  The global field,
host, character case and integer bounds have stayed fixed.
Thus the conditional G1 statement used by SP10 is completely
supplied, with its original coefficient-return mechanism.
\end{proof}

\paragraph{The resulting boundary.}
AX5 and AX6 are unconditional arithmetic suppliers relative to
the explicitly listed classical facts and the reconstructed BR7.
AX7--AX9 are implications on any actual good datum, with no
image-density premise.  AX10 uses
\texttt{SymplecticFourImageDense} exactly once for (AX.34).
It does not supply the automorphic catalogue, actual model
construction, geometric point comparison or rigidity.  Those
are the separately recorded G2 and catalogue chains.  In
particular the source formulas that later produce
\texttt{SymplecticFourImageDense} from automorphic hypotheses
are continuation targets, not classical consequences of
the elementary construction above.
